\chapter{Diagonalization and Spectral Decomposition}
\label{ch:diagonalization-spectral-decomposition}

\begin{draftbox}
This chapter is scaffolded. It will absorb the diagonalization, similarity, and
partial diagonalization material.
\end{draftbox}

\section{Guiding Question}

When can a matrix be made diagonal by changing coordinates?

\section{Planned Core Ideas}

\begin{itemize}
  \item an eigenbasis diagonalizes a linear map;
  \item $A=P\Lambda P^{-1}$ turns powers into scalar powers;
  \item similar matrices represent the same map in different coordinates;
  \item repeated eigenvalues can fail to provide enough eigenvectors;
  \item partial diagonalization still exposes invariant structure.
\end{itemize}

\section{Planned Python Thread}

Diagonalize small matrices, compute powers using $P\Lambda^kP^{-1}$, and test
failure cases.

\section*{Chapter Summary}
\addcontentsline{toc}{section}{Chapter Summary}

To be completed in the chapter-shaping pass.

\section*{Exercises}
\addcontentsline{toc}{section}{Exercises}

\subsection*{A. Theory}

\begin{exercise}
\exerciseid{13.A1}
State the spectral theorem for real symmetric matrices and explain why
orthogonality is part of the conclusion.
\end{exercise}

\begin{exercise}
\exerciseid{13.A2}
Let $A$ be symmetric. Prove that eigenvectors with distinct eigenvalues are
orthogonal.
\end{exercise}

\subsection*{B. By Hand}

\begin{exercise}
\exerciseid{13.B1}
Compute the Rayleigh quotient of
\[
  A=\begin{bmatrix}2&0\\0&5\end{bmatrix}
\]
at $x=(1,1)^T$ and at $x=(0,1)^T$.
\end{exercise}

\begin{exercise}
\exerciseid{13.B2}
For the same matrix $A$, find the minimum and maximum possible values of the
Rayleigh quotient.
\end{exercise}

\subsection*{C. Python / NumPy}

\begin{exercise}
\exerciseid{13.C1}
Sample random unit vectors and evaluate the Rayleigh quotient for a symmetric
$2\times 2$ matrix. Compare the observed minimum and maximum with
\texttt{np.linalg.eigvalsh}.
\end{exercise}

\begin{exercise}
\exerciseid{13.C2}
Implement 20 steps of power iteration for a symmetric matrix. Compare the final
Rayleigh quotient with the eigenvalue of largest absolute value.
\end{exercise}

\subsection*{D. Challenge}

\begin{exercise}
\exerciseid{13.D1}
Explain why the maximum Rayleigh quotient of a symmetric matrix should be an
eigenvalue. Use the orthonormal eigenbasis picture.
\end{exercise}

